%% kapa: exam=1 round=36 year=2025 period=2 form=A subject=appraisal_law subject_ko=감정평가관계법규 seq=3 number=3
%% source: Q-Net 감정평가사 기출문제 (cst00309) artlSeq=5231525 file=raw/1008/5231525/2025년 제36회 감정평가사 1차시험 2교시 문제지 원본.pdf
%% conversion: kordoc 4.24.51 -> tools/gen1.py (auto); review=proofread-sonnet; flags=none
\begin{question}{3}
국토의 계획 및 이용에 관한 법령상 도시ㆍ군관리계획에 관한 설명으로 옳은 것은?
\begin{choices}
  \item 도시ㆍ군관리계획은 광역도시계획과 도시ㆍ군기본계획(생활권계획 제외)에 부합되어야 한다.
  \item 주민은 산업ㆍ유통개발진흥지구의 지정 및 변경에 관한 사항에 대하여 도시ㆍ군관리계획의 입안을 제안할 수 있다.
  \item 시ㆍ도지사는 개발제한구역의 지정 및 변경에 관한 도시ㆍ군관리계획을 직접 결정하여야 한다.
  \item 도시ㆍ군관리계획 결정의 효력은 지형도면을 고시한 날의 다음날부터 발생한다.
  \item 시장 또는 군수가 입안한 지구단위계획의 수립ㆍ변경에 관한 도시ㆍ군관리계획은 시ㆍ도지사가 결정한다.
\end{choices}
\end{question}
\answer{2}
